% !PRINTED 614
\weberpage{614}
%<<CONT
Let $\mathfrak{m}$ be an integral ideal in $\Omega$, and let $A$ consist of all integral and fractional numbers $\alpha$ in $\Omega$ that are coprime to $\mathfrak{m}$ and satisfy the condition:
\weq{1}{\alpha \equiv 1 \quad (\mathrm{mod}\ \mathfrak{m}).}

If we take the prime divisors of $\mathfrak{m}$ into the excludent, and denote by $O$ the group of the numbers in $\Omega$, then $O$ splits with respect to the modulus $\mathfrak{m}$ into $\psi(\mathfrak{m})$ classes of numbers, in which $\psi(\mathfrak{m})$ has the meaning it has in Bd.\ II, \S\,168, namely, when $\mathfrak{p}$ runs through the prime divisors of $\mathfrak{m}$:
\weq{2}{\psi(\mathfrak{m}) = N(\mathfrak{m})\ \Pi\Bigl(1 - \frac{1}{N(\mathfrak{p})}\Bigr).}

To the group $A$ there belong not merely individual numbers, but classes of numbers with respect to the modulus $\mathfrak{m}$. Accordingly:
\weq{3}{(O, A) = \psi(\mathfrak{m}).}

In order to determine the class number for our group, we apply the theorems of \S\,161.

According to these, if $\overline{E}$ denotes the group of the functional units and $E$ the group of the numerical units:
\[
(\overline{O}, \overline{E}A) = (\overline{O}, \overline{E}O)\,(\overline{E}O, \overline{E}A),
\]
\[
(\overline{E}O, \overline{E}A) = (\overline{E}O, \overline{E}EA) = (O, EA).
\]

Further
\[
(O, EA)\,(EA, A) = (O, A) = \psi(\mathfrak{m}).
\]
$(EA, A) = e$ is the number of the units incongruent with respect to $\mathfrak{m}$, and
\[
(\overline{O}, \overline{E}O) = H
\]
is the class number of the field $\Omega$,
\[
(\overline{O}, \overline{E}A) = h
\]
is the class number with respect to $A$. From this there results:
\weq{4}{h = \frac{H\,\psi(\mathfrak{m})}{e}.}

The number $e$ is in general equal to the number of the units in $\Omega$, hence $= 2$, and in the two exceptional cases $\Delta = -4$, $\Delta = -3$ we have $e = 4$ and $e = 6$. The number $e$ is smaller when among the units there occur congruent ones (modulo $\mathfrak{m}$), which is only possible when $\mathfrak{m}$ is a divisor of 2 or of 3.
